\section*{Anexo H. Segundo estudio (paso 3): llegadas aleatorias y un fijo ajustado al segundo}
\phantomsection
\addcontentsline{toc}{subsection}{Anexo H. Segundo estudio (paso 3): llegadas aleatorias y un fijo ajustado al segundo}
\setcounter{table}{0}
\renewcommand{\thetable}{H.\arabic{table}}
\renewcommand{\theHtable}{H.\arabic{table}}
\setcounter{figure}{0}
\renewcommand{\thefigure}{H.\arabic{figure}}
\renewcommand{\theHfigure}{H.\arabic{figure}}
\setcounter{equation}{0}
\renewcommand{\theequation}{H.\arabic{equation}}
\renewcommand{\theHequation}{H.\arabic{equation}}
% Segundo estudio (plan 3, 05-09.11.2026). Registro previo: plan3/PRERREGISTRO.md (tag plan3-prerregistro-A, 05.11.2026
% 11:58), enmiendas plan3/doc/enmienda_1.md (06.11) y enmienda_2.md (07.11), parte B plan3/doc/parte_B.md (07.11),
% cierre fc74335 (09.11.2026 17:56) y veredicto 2eb76ff. Todas las cifras salen de plan3/salidas/ (veredicto3.md y sus
% CSV, escritos por plan3/veredicto3.py --final); las tablas plan3_*.tex y la Figura fig:plan3-cambio las escribe
% tablas_plan3.py. Nombres del glosario: "veredicto" (en el registro, "etiqueta"), "IPPO con mensajes corregido" (V2),
% "fijo ajustado al segundo" (fijo fino), "llegadas regulares" y "llegadas aleatorias".

Este anexo da el detalle del segundo estudio (paso 3), con el que la
Sección~\ref{sec:resultado-central} responde si el modelo reduce el
retraso total con llegadas aleatorias y frente a un fijo ajustado al
segundo. El diseño está en la Sección~\ref{sec:diseno-segundo}.

% fuente: plan3/salidas/veredicto3.md ("Respuesta corta": numero unico -10.8 % [-11.4, -10.2] y -11.1 % [-12.0, -10.4]
% con llegadas aleatorias; -0.4 % [-1.0, 0.1] y -0.4 % [-0.8, 0.1] con regulares; H5 pasa en 3 de 6: P1, P2 y P6) y
% tabla "Referencias al lado" (P1, P2 y P6: MEJORA o FUNCIONA frente a todas); veredicto3_celdas.csv (P3 +1.6 y +1.6;
% P4 -0.5 y 0.0; P5 +0.2 y +0.3)

\textbf{Cómo se lee.} Los apartados van en este orden: las reglas del
registro previo y lo que cambia frente al primer estudio, el actuado
ajustado de cada caso, el resultado con llegadas aleatorias y con
regulares, el resumen en un número, las otras referencias, el modelo
citado, las hipótesis, la sensibilidad, con la mezcla de vehículos de
Lima, el margen de mejora de cada caso, lo esperado y las desviaciones. Un cambio negativo es menos retraso que la
referencia. Las semillas de cada etapa están en la
Tabla~\ref{tab:plan2-semillas}; el ajuste de las referencias, en la
Sección~\ref{sec:referencias-segundo}; los modelos, en la
Sección~\ref{sec:aprendices}, y los programas, en la
Tabla~\ref{tab:programas-segundo}.

\subsection*{Registro previo}

El segundo estudio es posterior al primero y no cambia ninguno de sus
resultados: tiene sus propias reglas, publicadas en el repositorio
antes de entrenar. El registro previo fija los casos, los modelos, las
referencias, las semillas, la regla de veredicto, las hipótesis y lo
que se esperaba, y ningún criterio se cambió después de ver
resultados. Lo completan tres resguardos y dos enmiendas:
\begin{itemize}
  \item \textbf{Referencias congeladas.} Las referencias ajustadas se
    publicaron antes de evaluar.
  \item \textbf{Evaluación a ciegas.} Los resultados de las tandas se
    guardaron sin abrir hasta el cierre del estudio.
  \item \textbf{Un solo cálculo del veredicto.} Lo calculó un programa,
    una sola vez, al cierre.
  \item \textbf{Enmienda 1.} El estudio se centra en las llegadas
    aleatorias; con llegadas regulares se quitan casi todas las
    variantes.
  \item \textbf{Enmienda 2.} Con llegadas regulares queda solo IPPO con
    mensajes corregido frente al fijo ajustado al segundo, con el
    actuado ajustado, en P1 a P6; se quita la política única por red.
\end{itemize}
Las dos enmiendas reducen el alcance y se publicaron antes de evaluar
ninguna tanda; las hipótesis que retiran se listan junto a la
Tabla~\ref{tab:plan3-hipotesis}.
% fuente: plan3/PRERREGISTRO.md §0 y §14; tags plan3-prerregistro-A (05.11.2026), plan3-enmienda-1 (06.11),
% plan3-enmienda-2 y plan3-parte-B (07.11); cierre fc74335 (09.11.2026 17:56); plan3/cola/trabajos.jsonl (1944 trabajos;
% 201 cadenas de entrenamiento, 44 160 episodios; 642 + 98 evaluaciones del examen) y cola/estado.json (1002 h-proceso)
% fuente de las enmiendas: plan3/doc/enmienda_1.md y enmienda_2.md; plan3/doc/desviaciones.md, §4 y §5 (revision
% final del 22.11.2026: las enmiendas pasan aqui, sin fechas, desde "Enmiendas y desviaciones")

\subsection*{Qué cambia frente al primer estudio}

La Tabla~\ref{tab:plan3-cambia} compara los dos estudios. La red, la
demanda media de cada calle, la recompensa y el algoritmo son los del
primero. Con llegadas aleatorias cambia el instante en que llega cada
vehículo; el fijo y el actuado se vuelven a ajustar con esas llegadas,
y el modelo se entrena con ellas, siempre con semillas distintas de las
de evaluación. El modelo citado de cada caso se elige con la regla de
la Sección~\ref{sec:diseno-segundo} (Tabla~\ref{tab:plan3-promocion}).
% fuente: plan3/PRERREGISTRO.md §0 (modelo base y modelo citado) y §12 (regla de promocion; el modelo completo lleva
% la observacion nueva de 75 m, las tres limpiezas, 400 episodios con piso de la tasa y el reloj de 2 s);
% plan3/generar_trabajos3.py, OP_COMPLETO; anexo_g_diagnostico.tex, "Las dos correcciones"

\subsection*{El actuado ajustado de cada caso}

El actuado ajustado parte del actuado del primer estudio, con los
valores por defecto de SUMO, y prueba un parámetro a la vez, en este
orden y con las semillas 2001 a 2003: el hueco máximo entre vehículos
que sostiene el verde (1, 1.5, 2, 3 y 4~s), la distancia del detector,
en segundos de viaje a la velocidad máxima (1, 2, 3 y 4~s), el umbral
de atasco (ninguno, 10 y 30~s) y el verde máximo de la avenida (30, 45
y 60~s) y de la transversal (20, 30 y 60~s). Cada paso parte del mejor
de los anteriores. Después prueba la variante de SUMO que decide por el
retraso de los vehículos, con un alcance del detector de 50, 100 o
200~m y un retraso mínimo de 1 o 3~s, y las dos variantes coordinadas
con el ciclo y los desfases del fijo ajustado al segundo. Al final, las
tres mejores se vuelven a simular con las semillas 2001 a 2010 (con
llegadas aleatorias, las diez mejores con las 2001 a 2030), y gana la
de menor retraso total medio. La Tabla~\ref{tab:plan3-actuado} da la
elegida en cada caso.
% fuente: docstring de plan3/actuado.py (pasos 1 a 6 y final: rejillas 1/1.5/2/3*/4, 1/2*/3/4, -1*/10/30, 30/45/60*
% y 20/30/60*; delay_based con detectorRange 50/100/200 x minTimeLoss 1*/3; coordinadas (a) y (b) "fin fijo"; final
% con 2001-2010, o las 10 mejores con 2001-2030 con llegadas aleatorias; * = defecto de SUMO, el punto de partida);
% semillas_pasos y semillas de cada plan3/barridos/<esc>/actuado_aj.json (22 variantes probadas en cada caso)

% fuente: plan3/barridos/<esc>/actuado_aj.json, campo "variante" (tipo, max-gap, detector-gap, jam-threshold,
% maxdur_av, maxdur_tr, coordinado y fin_fijo; delay_based: detectorRange y minTimeLoss), con <esc> = poi_<caso> para
% llegadas aleatorias y <caso> para regulares (P6 con la entrada corregida: malla3_alta_e); el programa congelado es
% plan3/<esc>_actuado_aj.add.xml (sha256 en plan3/referencias_congeladas.json). Los de la sensibilidad (mix_corredor,
% mix_red, per_corredor y per_red) estan en los mismos archivos y no se tabulan
\begin{table}[H]
\setstretch{1}
\centering
\footnotesize
\setlength{\tabcolsep}{3pt}
\begin{tabular}{>{\raggedright\arraybackslash}p{2.6cm}lcccc>{\raggedright\arraybackslash}p{2.0cm}}
\toprule
Caso & Variante & \makecell{Hueco\\máximo (s)} & \makecell{Distancia del\\detector (s)} & \makecell{Umbral de\\atasco (s)} & \makecell{Verde máximo,\\avenida /\\transversal (s)} & Coordinado \\
\midrule
\multicolumn{7}{l}{\textit{Llegadas aleatorias}} \\
P1 & por retraso & \multicolumn{3}{c}{alcance 200~m, retraso 1~s} & 60 / 20 & no \\
\specialrule{0.25pt}{1.5pt}{1.5pt}
P2 & actuado & 1.5 & 2 & ninguno & 60 / 60 & sí, con fin fijo \\
\specialrule{0.25pt}{1.5pt}{1.5pt}
P3 & actuado & 1.5 & 2 & ninguno & 60 / 60 & sí, con fin fijo \\
\specialrule{0.25pt}{1.5pt}{1.5pt}
P4 & actuado & 1.5 & 2 & ninguno & 60 / 60 & sí, con fin fijo \\
\specialrule{0.25pt}{1.5pt}{1.5pt}
P5 & actuado & 1.5 & 2 & ninguno & 60 / 60 & sí \\
\specialrule{0.25pt}{1.5pt}{1.5pt}
P6 & actuado & 1.5 & 2 & 30 & 60 / 60 & no \\
\midrule
\multicolumn{7}{l}{\textit{Llegadas regulares}} \\
P1 & actuado & 3 & 3 & ninguno & 60 / 60 & sí \\
\specialrule{0.25pt}{1.5pt}{1.5pt}
P2 & actuado & 3 & 4 & 30 & 60 / 60 & sí, con fin fijo \\
\specialrule{0.25pt}{1.5pt}{1.5pt}
P3 & actuado & 2 & 2 & ninguno & 60 / 60 & sí, con fin fijo \\
\specialrule{0.25pt}{1.5pt}{1.5pt}
P4 & actuado & 1 & 2 & ninguno & 60 / 60 & sí, con fin fijo \\
\specialrule{0.25pt}{1.5pt}{1.5pt}
P5 & actuado & 1 & 2 & ninguno & 60 / 60 & sí \\
\specialrule{0.25pt}{1.5pt}{1.5pt}
P6 & actuado & 1 & 1 & 10 & 60 / 60 & sí, con fin fijo \\
\specialrule{0.25pt}{1.5pt}{1.5pt}
P6 con la entrada corregida & actuado & 1.5 & 2 & 10 & 30 / 60 & sí, con fin fijo \\
\bottomrule
\end{tabular}
\caption[Segundo estudio: actuado ajustado de cada caso]{Actuado ajustado de cada caso: la variante elegida y sus parámetros. Actuado: el actuado de SUMO; por retraso: su variante que decide por el retraso de los vehículos que se acercan, con el alcance del detector y el retraso mínimo que cuenta. Coordinado: con el ciclo y los desfases del fijo ajustado al segundo (Tabla~\ref{tab:fijo-segundo}), y el verde de la avenida sin terminar antes que en el plan fijo; con fin fijo, además, la avenida termina siempre donde el plan fijo y la transversal nunca después. El verde mínimo es de 10~s.}
\label{tab:plan3-actuado}
\end{table}

\subsection*{Resultado con llegadas aleatorias}

La Tabla~\ref{tab:plan3-aleatorias} da el veredicto de cada caso:
FUNCIONA en P1, P2 y P6 (en P2, con el entrenamiento de su versión
evaluada sin estabilizar), FALLA en P3, por un 1.6\,\%, y EMPATA en P4
y P5.

\begin{table}[H]
\setstretch{1}
\centering
\footnotesize
\setlength{\tabcolsep}{3pt}
% Generado por tablas_plan3.py a partir de plan3/salidas/. No editar a mano.
% Segundo estudio, llegadas aleatorias: etiqueta citable del modelo citado frente al fijo ajustado al
% segundo, con el cambio del retraso total y su IC95 en cada tanda (100 semillas por tanda).
% fuente: plan3/salidas/veredicto3_celdas.csv (filas P1-P6 con llegadas aleatorias: modelo_citado,
% etiqueta_citable, cambio_T1, ic_T1, cambio_T2, ic_T2, repite_T1 y repite_T2).
% Entorno: \setstretch{1}, \footnotesize y \setlength{\tabcolsep}{3pt}.
\begin{tabular}{>{\raggedright\arraybackslash}p{2.9cm}>{\raggedright\arraybackslash}p{2.2cm}>{\raggedright\arraybackslash}p{2.2cm}>{\raggedright\arraybackslash}p{2.4cm}>{\raggedright\arraybackslash}p{2.4cm}>{\raggedright\arraybackslash}p{1.3cm}}
\toprule
Caso & Modelo citado & Veredicto & Cambio, tanda 1 (\%) & Cambio, tanda 2 (\%) & Se repite \\
\midrule
P1, corredor, demanda nominal (1810~veh/h) & Modelo completo & FUNCIONA & $-16.3$ [$-17.7$, $-15.0$] & $-17.1$ [$-18.6$, $-15.8$] & 3 de 3 y 3 de 3 \\
\specialrule{0.25pt}{1.5pt}{1.5pt}
P2, corredor, demanda alta (3620~veh/h) & Modelo completo & FUNCIONA (sin estabilizar) & $-19.8$ [$-20.5$, $-19.0$] & $-20.2$ [$-20.9$, $-19.6$] & 3 de 3 y 3 de 3 \\
\specialrule{0.25pt}{1.5pt}{1.5pt}
P3, red con rotonda, demanda nominal (2150~veh/h) & IPPO con mensajes corregido & FALLA & $+1.6$ [$+1.4$, $+1.7$] & $+1.6$ [$+1.4$, $+1.8$] & 5 de 5 y 5 de 5 \\
\specialrule{0.25pt}{1.5pt}{1.5pt}
P4, red con rotonda, demanda alta (2580~veh/h) & Modelo completo & EMPATA & $-0.5$ [$-1.3$, $+0.2$] & $0.0$ [$-0.7$, $+0.6$] & 3 de 3 y 3 de 3 \\
\specialrule{0.25pt}{1.5pt}{1.5pt}
P5, malla, demanda nominal (4350~veh/h) & IPPO con mensajes corregido & EMPATA & $+0.2$ [$-0.1$, $+0.5$] & $+0.3$ [$0.0$, $+0.7$] & 0 de 5 y 0 de 5 \\
\specialrule{0.25pt}{1.5pt}{1.5pt}
P6, malla, demanda alta (8700~veh/h) & Modelo completo & FUNCIONA & $-25.7$ [$-28.4$, $-23.0$] & $-26.9$ [$-30.8$, $-23.5$] & 3 de 3 y 3 de 3 \\
\bottomrule
\end{tabular}

\caption[Segundo estudio: resultado con llegadas aleatorias]{Resultado con llegadas aleatorias: cambio del retraso total del modelo citado frente al fijo ajustado al segundo, con su IC95 entre corchetes, en cada tanda de 100 semillas. Veredicto: el que se repite en las dos tandas, con la regla de la Tabla~\ref{tab:veredictos} y el actuado ajustado en el lugar del actuado. Se repite: en cuántos entrenamientos del modelo, evaluados uno por uno, el cambio queda por debajo de cero (por encima, si el veredicto es FALLA), en la tanda 1 y en la tanda 2.}
\label{tab:plan3-aleatorias}
\end{table}

La lectura de cada caso está en la Sección~\ref{sec:resultado-casos}.
El registro pide además dos lecturas: en P3 el modelo reproduce el
ciclo mínimo, con su mismo retraso, y en P3, P4 y P5 no supera a la
regla simple, mientras que el fijo de verdes de 5 en 5 lo iguala o lo
supera (Tabla~\ref{tab:plan3-referencias}). Sin los cambios del modelo
completo, IPPO con mensajes corregido recibe MEJORA en P1 ($-6.3$ y
$-8.7$\,\%), FUNCIONA en P2 ($-5.7$ y $-5.8$\,\%), con un entrenamiento
que sí se estabilizó, y en P6 ($-18.4$ y $-20.2$\,\%), y FALLA en P4
($+1.2$ y $+1.5$\,\%), donde el modelo completo empata.
% fuente: plan3/salidas/veredicto3_celdas.csv (filas con llegadas aleatorias) y veredicto3.md, "Lecturas obligadas";
% IPPO con mensajes corregido frente al fijo ajustado al segundo: veredicto3_comparaciones.csv, clave v2poi_version,
% referencia "fijo fino" (P1 -6.32 y -8.70 MEJORA; P2 -5.65 y -5.78 FUNCIONA; P4 +1.24 y +1.48 FALLA; P6 -18.37 y
% -20.24 FUNCIONA); P3: igual_ciclo_minimo = si

\subsection*{Resultado con llegadas regulares}

Con llegadas regulares solo se evalúa IPPO con mensajes corregido
(Tabla~\ref{tab:plan3-regulares}; lectura en la
Sección~\ref{sec:resultado-casos}). El fijo ajustado al segundo usa
verdes que el modelo, al decidir cada 5~s, no puede dar. Frente al
actuado ajustado, el modelo mejora en P2, empata en P4 y P6 y tiene más
retraso en P1, P3 y P5. En P6 la mitad del retraso total es espera para
entrar a la red, igual con todos los controles, y diluye el cambio; por
eso P6 se evalúa también con la entrada corregida, en la que el modelo
tiene un 10\,\% más de retraso. En P1, la versión evaluada de IPPO con
mensajes corregido en el diagnóstico del corredor (Anexo~G), que sale de
otro entrenamiento, da $-2.7$\,\% en las dos tandas frente al fijo
ajustado al segundo, también MEJORA.

\begin{table}[H]
\setstretch{1}
\centering
\footnotesize
\setlength{\tabcolsep}{3pt}
% Generado por tablas_plan3.py a partir de plan3/salidas/. No editar a mano.
% Segundo estudio, llegadas regulares: etiqueta citable de IPPO con mensajes corregido frente al fijo
% ajustado al segundo (30 semillas por tanda).
% fuente: plan3/salidas/veredicto3_celdas.csv (filas P1-P6 y P6e con llegadas regulares).
% Entorno: \setstretch{1}, \footnotesize y \setlength{\tabcolsep}{3pt}.
\begin{tabular}{>{\raggedright\arraybackslash}p{4.0cm}>{\raggedright\arraybackslash}p{3.0cm}>{\raggedright\arraybackslash}p{2.5cm}>{\raggedright\arraybackslash}p{2.5cm}>{\raggedright\arraybackslash}p{1.4cm}}
\toprule
Caso & Veredicto & Cambio, tanda 1 (\%) & Cambio, tanda 2 (\%) & Se repite \\
\midrule
P1, corredor, demanda nominal (1810~veh/h) & MEJORA & $-2.8$ [$-3.7$, $-2.0$] & $-2.8$ [$-3.6$, $-1.9$] & 5 de 5 y 5 de 5 \\
\specialrule{0.25pt}{1.5pt}{1.5pt}
P2, corredor, demanda alta (3620~veh/h) & EMPATA & $-1.7$ [$-3.3$, $0.0$] & $-1.3$ [$-3.0$, $+0.4$] & 4 de 5 y 5 de 5 \\
\specialrule{0.25pt}{1.5pt}{1.5pt}
P3, red con rotonda, demanda nominal (2150~veh/h) & FALLA & $+1.7$ [$+1.3$, $+2.1$] & $+1.7$ [$+1.3$, $+2.1$] & 5 de 5 y 5 de 5 \\
\specialrule{0.25pt}{1.5pt}{1.5pt}
P4, red con rotonda, demanda alta (2580~veh/h) & EMPATA & $+0.5$ [$-1.5$, $+2.5$] & $+0.4$ [$-1.1$, $+1.8$] & 0 de 5 y 0 de 5 \\
\specialrule{0.25pt}{1.5pt}{1.5pt}
P5, malla, demanda nominal (4350~veh/h) & EMPATA (sin estabilizar) & $+0.1$ [$-0.1$, $+0.5$] & $+0.1$ [$-0.2$, $+0.4$] & 0 de 5 y 1 de 5 \\
\specialrule{0.25pt}{1.5pt}{1.5pt}
P6, malla, demanda alta (8700~veh/h) & EMPATA & $-0.3$ [$-2.0$, $+1.8$] & $-0.2$ [$-1.4$, $+1.3$] & 4 de 5 y 2 de 5 \\
\specialrule{0.25pt}{1.5pt}{1.5pt}
P6 con la entrada corregida & FALLA & $+9.8$ [$+9.2$, $+10.5$] & $+10.5$ [$+9.6$, $+11.5$] & 5 de 5 y 5 de 5 \\
\bottomrule
\end{tabular}

\caption[Segundo estudio: resultado con llegadas regulares]{Resultado con llegadas regulares: cambio del retraso total de IPPO con mensajes corregido frente al fijo ajustado al segundo, con su IC95, en cada tanda de 30 semillas. Las columnas se leen como en la Tabla~\ref{tab:plan3-aleatorias}. En la última fila, cada vehículo entra por el carril que le conviene.}
\label{tab:plan3-regulares}
\end{table}
% fuente: plan3/salidas/veredicto3_celdas.csv (filas con llegadas regulares); PRERREGISTRO.md §4 (en P6, con la
% entrada por defecto, el 55 % del retraso total es espera para entrar) y §5 ("Del fijo fino frente a IPPO");
% veredicto3.md, "Referencias al lado" (regulares, frente al actuado ajustado: P1 FALLA, P2 MEJORA, P3 FALLA, P4
% EMPATA, P5 FALLA, P6 EMPATA) y "V2 del diagnostico del corredor" (P1: base 42 en el diagnostico, 16042 ahora);
% veredicto3_comparaciones.csv (P1, clave v2reg_s42: -2.67 y -2.68, MEJORA)

\subsection*{Resumen en un número}

Se calcula como en el primer estudio
(Ecuación~\ref{eq:media-geometrica}), sobre los seis casos:
\begin{itemize}
  \item \textbf{Llegadas aleatorias, frente al fijo ajustado al
    segundo:} $-10.8$\,\% [$-11.4$, $-10.2$] en la tanda 1 y
    $-11.1$\,\% [$-12.0$, $-10.4$] en la tanda 2. Incluye P2, cuyo
    entrenamiento quedó sin estabilizar.
  \item \textbf{Llegadas aleatorias, frente al fijo ajustado desde
    10~s:} $-15.8$\,\% [$-17.4$, $-14.2$] y $-14.5$\,\% [$-15.9$,
    $-13.1$].
  \item \textbf{Llegadas aleatorias, IPPO con mensajes corregido en
    lugar del modelo citado:} $-4.8$\,\% [$-5.6$, $-3.9$] y
    $-5.6$\,\% [$-6.5$, $-4.8$].
  \item \textbf{Llegadas regulares, frente al fijo ajustado al
    segundo:} $-0.4$\,\% [$-1.0$, $+0.1$] y $-0.4$\,\% [$-0.8$,
    $+0.1$]. Incluye P5, cuyo entrenamiento quedó sin estabilizar.
\end{itemize}
% fuente: plan3/salidas/veredicto3.md, "Respuesta corta" (numero unico por tipo de llegadas y referencia)

\subsection*{Frente a las otras referencias}

Además del fijo ajustado al segundo, cada caso tiene otras nueve
referencias. La Tabla~\ref{tab:plan3-referencias} da, para cada una, los
casos en que el modelo citado recibe cada veredicto: en P1, P2 y P6
tiene menos retraso que todas, y en P3, P4 y P5 no supera al actuado
ajustado, a la regla simple ni al ciclo mínimo.

\begin{table}[H]
\setstretch{1}
\centering
\footnotesize
\setlength{\tabcolsep}{3pt}
% Generado por tablas_plan3.py a partir de plan3/salidas/. No editar a mano.
% Segundo estudio, llegadas aleatorias: casos en que el modelo citado recibe cada etiqueta citable frente
% a cada una de las otras referencias (la misma regla, con esa referencia en el lugar del fijo ajustado al
% segundo).
% fuente: plan3/salidas/veredicto3.md, tabla 'Referencias al lado' (filas P1-P6 con llegadas aleatorias).
% a: mejora en una sola tanda. Entorno: \setstretch{1}, \footnotesize y \setlength{\tabcolsep}{3pt}.
\begin{tabular}{>{\raggedright\arraybackslash}p{4.2cm}>{\raggedright\arraybackslash}p{3.0cm}>{\raggedright\arraybackslash}p{1.8cm}>{\raggedright\arraybackslash}p{3.0cm}>{\raggedright\arraybackslash}p{1.8cm}}
\toprule
Frente a & FUNCIONA & MEJORA & EMPATA & FALLA \\
\midrule
Actuado ajustado & P2, P6 & P1 & P4, P5 & P3 \\
Actuado por defecto & P1 a P6 & --- & --- & --- \\
Fijo ajustado desde 10~s & P1, P2, P6 & P4 & P3\textsuperscript{a}, P5 & --- \\
Fijo de verdes de 5 en 5 & P1, P2, P6 & --- & P4, P5 & P3 \\
Fijo por franja & P1, P2, P6 & --- & P4, P5 & P3 \\
Regla simple & P1, P2, P6 & --- & P3, P4\textsuperscript{a}, P5 & --- \\
Regla reactiva & P1 a P6 & --- & --- & --- \\
Max-Pressure & P1 a P6 & --- & --- & --- \\
Ciclo mínimo & P1, P2, P6 & --- & P3, P4\textsuperscript{a}, P5 & --- \\
\bottomrule
\end{tabular}

\caption[Segundo estudio: el modelo citado frente a las otras referencias]{Casos en que el modelo citado recibe cada veredicto frente a cada una de las otras referencias, con llegadas aleatorias: la misma regla, con esa referencia en el lugar del fijo. Actuado por defecto: el actuado del plan. Fijo ajustado desde 10~s: el del Anexo~F. Fijo de verdes de 5 en 5: el mejor fijo con verdes múltiplos de 5~s. Fijo por franja: un plan fijo por cada media hora. Regla simple: ciclo mínimo, y alarga el verde de la avenida si hay cola. Regla reactiva: alarga el verde mientras llegan vehículos. (a) La mejora aparece en una sola tanda.}
\label{tab:plan3-referencias}
\end{table}

La Tabla~\ref{tab:plan3-referencias-cifras} da las cifras frente a dos
de ellas, el actuado ajustado y el fijo ajustado desde 10~s, en cada
caso y tanda.

\begin{table}[H]
\setstretch{1}
\centering
\footnotesize
\setlength{\tabcolsep}{3pt}
% Generado por tablas_plan3.py a partir de plan3/salidas/. No editar a mano.
% Segundo estudio, llegadas aleatorias: cambio del retraso total del modelo citado frente al actuado
% ajustado y frente al fijo entregado, en la tanda 1 y en la tanda 2 (%).
% fuente: plan3/salidas/veredicto3_comparaciones.csv (clave del modelo citado de cada caso, referencias
% 'actuado ajustado' y 'fijo entregado', metrica retraso_total_prom, columna cambio).
% Entorno: \setstretch{1}, \footnotesize y \setlength{\tabcolsep}{3pt}.
\begin{tabular}{>{\raggedright\arraybackslash}p{5.6cm}>{\raggedright\arraybackslash}p{3.6cm}>{\raggedright\arraybackslash}p{3.6cm}}
\toprule
Caso & Frente al actuado ajustado (\%) & Frente al fijo entregado (\%) \\
\midrule
P1, corredor, demanda nominal (1810~veh/h) & $-3.3$ y $-3.8$ & $-18.2$ y $-18.8$ \\
P2, corredor, demanda alta (3620~veh/h) & $-19.0$ y $-19.5$ & $-27.3$ y $-27.5$ \\
P3, red con rotonda, demanda nominal (2150~veh/h) & $+1.6$ y $+1.6$ & $-0.2$ y $-0.3$ \\
P4, red con rotonda, demanda alta (2580~veh/h) & $-0.5$ y $0.0$ & $-1.5$ y $-1.5$ \\
P5, malla, demanda nominal (4350~veh/h) & $0.0$ y $+0.3$ & $0.0$ y $+0.1$ \\
P6, malla, demanda alta (8700~veh/h) & $-11.5$ y $-12.9$ & $-39.0$ y $-32.4$ \\
\bottomrule
\end{tabular}

\caption[Segundo estudio: cambio frente al actuado ajustado y al fijo ajustado desde 10~s]{Cambio del retraso total del modelo citado frente al actuado ajustado y frente al fijo ajustado desde 10~s, con llegadas aleatorias, en la tanda 1 y en la tanda 2.}
\label{tab:plan3-referencias-cifras}
\end{table}
% fuente: plan3/salidas/veredicto3.md, "Referencias al lado", y veredicto3_comparaciones.csv (clave del modelo citado
% de cada caso frente a "actuado ajustado" y "fijo entregado"); PRERREGISTRO.md §5 (definicion de cada referencia)

\subsection*{Qué modelo se cita en cada caso}

El modelo citado se elige con una validación aparte, de 16 semillas que
no se usan para nada más (Tabla~\ref{tab:plan3-promocion}). El modelo
completo supera a IPPO con mensajes corregido por al menos un 1\,\% en
P1, P2, P4 y P6, y es el citado en esos casos; en P3 y P5 ninguna
variante llega a ese margen.

\begin{table}[H]
\setstretch{1}
\centering
\footnotesize
\setlength{\tabcolsep}{3pt}
% Generado por tablas_plan3.py a partir de plan3/salidas/. No editar a mano.
% Segundo estudio, regla de promocion con llegadas aleatorias: cambio del retraso total de cada variante
% frente a IPPO con mensajes corregido en las 16 semillas de promocion (15009-15024), y modelo citado.
% fuente: plan3/promocion/poi_<caso>.json (candidatos: cambio_vs_v2_pct) y
% plan3/salidas/veredicto3_celdas.csv (modelo_citado). Se promueve la variante de menor retraso si le
% gana a IPPO con mensajes corregido por al menos 1 %.
% Entorno: \setstretch{1}, \footnotesize y \setlength{\tabcolsep}{3pt}.
\begin{tabular}{l>{\raggedright\arraybackslash}p{2.0cm}>{\raggedright\arraybackslash}p{2.0cm}>{\raggedright\arraybackslash}p{2.0cm}>{\raggedright\arraybackslash}p{2.0cm}>{\raggedright\arraybackslash}p{3.6cm}}
\toprule
Caso & Modelo completo (\%) & Observación nueva (\%) & 400 episodios (\%) & Decisión cada 2~s (\%) & Modelo citado \\
\midrule
P1 & $-9.04$ & $-8.06$ & $-1.27$ & $+0.15$ & Modelo completo \\
P2 & $-13.61$ & $-9.83$ & $-1.48$ & $+1.71$ & Modelo completo \\
P3 & $-0.36$ & $+0.30$ & $0.00$ & --- & IPPO con mensajes corregido \\
P4 & $-1.24$ & $+0.08$ & $0.00$ & --- & Modelo completo \\
P5 & $+1.87$ & $+5.90$ & $0.00$ & --- & IPPO con mensajes corregido \\
P6 & $-8.79$ & $-6.89$ & $-5.24$ & $+3.67$ & Modelo completo \\
\bottomrule
\end{tabular}

\caption[Segundo estudio: elección del modelo citado]{Elección del modelo citado con llegadas aleatorias: cambio del retraso total de cada variante frente a IPPO con mensajes corregido, en las 16 semillas de la validación de promoción. Se cita la variante de menor retraso si mejora al menos un 1\,\%. La decisión cada 2~s sola se probó en tres casos.}
\label{tab:plan3-promocion}
\end{table}
% fuente: plan3/promocion/poi_<caso>.json y plan3/salidas/veredicto3.md, "Regla de promocion"

\subsection*{Las hipótesis registradas}

La Tabla~\ref{tab:plan3-hipotesis} da el resultado de las hipótesis que
se evaluaron. Una hipótesis sobre el método pasa en un caso solo si la
mejora es significativa en las dos tandas y se repite en los
entrenamientos.

\begin{table}[H]
\setstretch{1}
\centering
\footnotesize
\setlength{\tabcolsep}{3pt}
\begin{tabular}{>{\raggedright\arraybackslash}p{7.2cm}>{\raggedright\arraybackslash}p{6.4cm}}
\toprule
Hipótesis & Resultado \\
\midrule
Con llegadas aleatorias, IPPO con mensajes corregido mejora al fijo ajustado al segundo en al menos 3 de 6 casos & Pasa: 3 de 6 (P1, P2 y P6) \\
\specialrule{0.25pt}{1.5pt}{1.5pt}
La ventaja frente al fijo es mayor con llegadas aleatorias que con regulares & Pasa en 2 de 6 casos (P2 y P6) \\
\specialrule{0.25pt}{1.5pt}{1.5pt}
Los mensajes mejoran al modelo, con llegadas aleatorias & Pasa en 1 de 6 casos: P6 ($-5.2$ y $-5.3$\,\%). En P3 dan lo mismo \\
\specialrule{0.25pt}{1.5pt}{1.5pt}
Ver los vehículos cercanos mejora a IPPO con mensajes corregido & Pasa en 3 de 6 casos: P1, P2 y P6 (entre $-5.9$ y $-9.7$\,\%) \\
\specialrule{0.25pt}{1.5pt}{1.5pt}
Entrenar 400 episodios mejora la versión evaluada & Pasa en 1 de 6 casos: P6 ($-5.3$ y $-5.4$\,\%) \\
\specialrule{0.25pt}{1.5pt}{1.5pt}
Decidir cada 2~s, en vez de cada 5~s, mejora al modelo & No pasa en ninguno de los 3 casos probados (P1, P2 y P6) \\
\specialrule{0.25pt}{1.5pt}{1.5pt}
El modelo completo mejora a IPPO con mensajes corregido & Pasa en 4 de 6 casos: P1, P2, P4 y P6 (entre $-1.5$ y $-13.0$\,\%) \\
\specialrule{0.25pt}{1.5pt}{1.5pt}
Con llegadas regulares, IPPO con mensajes corregido mejora al fijo ajustado al segundo en P6 & No pasa: EMPATA \\
\specialrule{0.25pt}{1.5pt}{1.5pt}
Con llegadas aleatorias, algún control que reacciona mejora al fijo ajustado al segundo (descriptiva) & En 3 de 6 casos: P1, P2 y P6 \\
\bottomrule
\end{tabular}
\caption[Segundo estudio: hipótesis registradas]{Hipótesis registradas que se evaluaron y su resultado. Entre paréntesis, el cambio en la tanda 1 y en la tanda 2, o su rango en los casos que pasan.}
\label{tab:plan3-hipotesis}
\end{table}
% fuente: plan3/salidas/veredicto3.md, "Hipotesis" (H5, H6, H7, H8b, H9, H14, H15, H3 y H0), y veredicto3_hipotesis.csv
% (cambio_1 y cambio_2 por caso: H7 P6 -5.23 y -5.33, P3 0.00 y 0.00; H8b P1 -8.35 y -8.60, P2 -9.67 y -8.91, P6 -7.77
% y -5.90; H9 P6 -5.30 y -5.41; H14 P1 +0.02 y +0.39, P2 +0.77 y +1.48, P6 +4.73 y +4.35; H15 P1 -8.80 y -9.94, P2
% -13.00 y -12.97, P4 -1.76 y -1.48, P6 -8.94 y -8.34)

Las dos enmiendas retiran del registro las hipótesis que solo usaban
llegadas regulares: la comparación con el
entrenamiento del primer estudio, los mensajes, los vehículos cercanos,
los ajustes del entrenamiento, los 400 episodios, la decisión cada 2~s
y el modelo completo con esas llegadas, y una sola política para todos
los niveles de demanda de una red.
% fuente: plan3/salidas/veredicto3.md (H1, H2, H4, H7r, H8a, H8c, H12, H10 y H9, H14 y H15 regulares: "no se corre")

\subsection*{Sensibilidad: mezcla de vehículos y perfil de demanda}

El corredor y la red con demanda nominal se prueban además con dos
cambios en el tráfico. El primero es una mezcla de vehículos como la de
Lima (Tabla~\ref{tab:plan3-mezcla}): SUMO sortea el tipo de cada
vehículo con esas proporciones, y las tasas de llegada se multiplican
por 0.8378, el inverso de los 1.1936 vehículos equivalentes que vale en
promedio un vehículo de la mezcla, de modo que la demanda del caso se
conserva en vehículos equivalentes. La mezcla se evalúa también con un arranque más lento:
cada vehículo tarda 0.5~s más en arrancar después de detenerse, solo al
evaluar. El segundo cambio es una demanda que sube y baja: la del caso
multiplicada por 2/3, 1, 4/3, 4/3, 1 y 2/3 en bloques de 10~minutos, con
el mismo total en cada media hora. Los dos los escribe el programa de
los escenarios del estudio (Tabla~\ref{tab:programas-segundo}).
% fuente: plan3/demanda_estocastica.py (JICA_2013: auto 541 088, moto 29 919 y publico 114 537 veh/dia, la suma por
% filas de la Tabla 2.16 sin camiones; ATU_2021: combi 6 420, micro 7 533 y bus 8 357; EQUIVALENTE; TIPOS y COMUNES;
% composicion() = 0.789283, 0.043643, 0.048078, 0.056413 y 0.062584; equivalencia_media() = 1.193600; PERFIL =
% (1.0, 1.5, 2.0, 2.0, 1.5, 1.0) / 1.5 en bloques de 600 s), igual en plan3/mix_corredor.rou.xml ("tasas x 0.837802 =
% 1/1.193600") y plan3/mixa_corredor.rou.xml (startupDelay 0.5 en los cinco tipos); escenarios_plan3.RECETAS (mix_,
% mixa_ con arranque 0.5 "solo para evaluar", y per_, en corredor y red). Fuentes revisadas en sus PDF el 22.11.2026:
% JICA 2013, Tabla 2.16, p. 2-21 (la fila Total publicada, 27 574 / 518 553 / 107 011, omite el puente Huachipa); Lima
% Como Vamos, reporte 2021, p. 16 (flota regular por antiguedad, 22 310 en total); UNFV, Tabla N.o 2, p. 28. Moto y
% bus: valores de su clase en SUMO 1.25, medidos con TraCI segun el comentario del programa; combi y micro: supuestos
% del programa. Casos de Lima: nudo_ovalo_estudio2/demo_rl/lima/demanda.py ("Las tasas son vehiculos reales por hora:
% no se convierten a vehiculos equivalentes"; CONDUCTA vacia) y corredor_estudio2/demo_rl/lima/corredor/corredor.py
% (TIPOS, COMUNES y composicion(), tasas sin escalar)
\begin{table}[H]
\setstretch{1}
\centering
\footnotesize
\setlength{\tabcolsep}{4pt}
\begin{tabular}{lrrrrr}
\toprule
\makecell[l]{Tipo\\(clase de SUMO)} & \makecell[r]{Proporción\\(\%)} & \makecell[r]{Largo\\(m)} & \makecell[r]{Aceleración\\(m/s$^2$)} & \makecell[r]{Deceleración\\(m/s$^2$)} & \makecell[r]{Vehículos\\equivalentes} \\
\midrule
Auto (passenger) & 78.93 & 5 & 2.6 & 4.5 & 1.00 \\
\specialrule{0.25pt}{1.5pt}{1.5pt}
Moto (motorcycle) & 4.36 & 2.2 & 6 & 10 & 1.00 \\
\specialrule{0.25pt}{1.5pt}{1.5pt}
Combi (bus) & 4.81 & 6 & 2 & 4.5 & 1.25 \\
\specialrule{0.25pt}{1.5pt}{1.5pt}
Microbús (bus) & 5.64 & 8 & 1.5 & 4 & 2.00 \\
\specialrule{0.25pt}{1.5pt}{1.5pt}
Bus (bus) & 6.26 & 12 & 1.2 & 4 & 3.00 \\
\bottomrule
\end{tabular}
\caption[Segundo estudio: mezcla de vehículos de Lima]{Mezcla de vehículos de Lima. Proporción: el tránsito diario por tipo en la línea cortina del río Rímac, sin camiones~\citep{JICA2013_Encuesta}, con el transporte público repartido entre combis, microbuses y buses como la flota regular de la ATU en 2021~\citep{LimaComoVamos2022_Reporte2021}. Largo, aceleración y deceleración: los del auto de los demás casos; para la moto y el bus, los de su clase en SUMO 1.25; para la combi y el microbús, supuestos. Vehículos equivalentes: factores del INVERMET para el auto, la camioneta rural (combi), el microbús y el ómnibus (bus)~\citep{UNFV2018_ImpactoVial}; la moto no tiene factor y se supone 1.00. Todos los tipos tienen $\sigma = 0.5$, 2.5~m de separación al detenerse y una velocidad máxima de 13.89~m/s (50~km/h); los demás parámetros son los de su clase en SUMO 1.25. Los casos de Lima usan la misma mezcla, con vehículos reales por hora, sin la conversión a vehículos equivalentes.}
\label{tab:plan3-mezcla}
\end{table}

La Tabla~\ref{tab:plan3-sensibilidad} da el resultado; su lectura está
en la Sección~\ref{sec:resultado-casos}. En la red, el modelo sigue con
menos de un 2\,\% más de retraso que el fijo, como con el tráfico de un
solo tipo de vehículo. Frente a las otras referencias, en el corredor
con la mezcla tiene un 7\,\% más de retraso que el actuado ajustado y no
supera a la regla reactiva; con el perfil empata con el actuado ajustado
y supera, por un 12\,\%, a un plan fijo que cambia cada 10 minutos.

\begin{table}[H]
\setstretch{1}
\centering
\footnotesize
\setlength{\tabcolsep}{3pt}
% Generado por tablas_plan3.py a partir de plan3/salidas/. No editar a mano.
% Segundo estudio, sensibilidad con llegadas aleatorias: mezcla de vehiculos de Lima, arranque lento y
% perfil de 10 min, en el corredor y en la red con demanda nominal.
% fuente: plan3/salidas/veredicto3_celdas.csv (filas mix1, mixa1, per1, mix3, mixa3 y per3).
% Entorno: \setstretch{1}, \footnotesize y \setlength{\tabcolsep}{3pt}.
\begin{tabular}{>{\raggedright\arraybackslash}p{4.4cm}>{\raggedright\arraybackslash}p{2.8cm}>{\raggedright\arraybackslash}p{2.4cm}>{\raggedright\arraybackslash}p{2.4cm}>{\raggedright\arraybackslash}p{1.4cm}}
\toprule
Caso & Veredicto & Cambio, tanda 1 (\%) & Cambio, tanda 2 (\%) & Se repite \\
\midrule
Corredor nominal, mezcla de vehículos de Lima & MEJORA & $-3.8$ [$-4.7$, $-2.9$] & $-3.2$ [$-4.1$, $-2.4$] & 3 de 3 y 3 de 3 \\
Corredor nominal, mezcla de Lima y arranque lento & MEJORA & $-3.4$ [$-4.3$, $-2.5$] & $-3.2$ [$-4.0$, $-2.3$] & 3 de 3 y 3 de 3 \\
Corredor nominal, demanda con perfil de 10~min & FUNCIONA (sin estabilizar) & $-12.1$ [$-14.5$, $-9.7$] & $-9.8$ [$-11.8$, $-7.8$] & 3 de 3 y 3 de 3 \\
\midrule
Red nominal, mezcla de vehículos de Lima & FALLA & $+1.7$ [$+1.5$, $+1.9$] & $+1.8$ [$+1.7$, $+2.0$] & 3 de 3 y 3 de 3 \\
Red nominal, mezcla de Lima y arranque lento & FALLA & $+1.3$ [$+1.1$, $+1.4$] & $+1.2$ [$+1.0$, $+1.3$] & 3 de 3 y 3 de 3 \\
Red nominal, demanda con perfil de 10~min & FALLA & $+1.0$ [$+0.5$, $+1.5$] & $+1.1$ [$+0.6$, $+1.7$] & 3 de 3 y 3 de 3 \\
\bottomrule
\end{tabular}

\caption[Segundo estudio: sensibilidad a la mezcla de vehículos y al perfil de demanda]{Sensibilidad con llegadas aleatorias: cambio del retraso total de IPPO con mensajes corregido, entrenado con cada tráfico, frente al fijo ajustado al segundo con ese tráfico. El arranque lento solo se usa al evaluar. Las columnas se leen como en la Tabla~\ref{tab:plan3-aleatorias}.}
\label{tab:plan3-sensibilidad}
\end{table}
% fuente: plan3/salidas/veredicto3_celdas.csv (mix1, mixa1, per1, mix3, mixa3 y per3); plan3/demanda_estocastica.py
% (mezcla: JICA 2013 sin camiones y flota de la ATU de 2021; perfil de 10 min con factores 2/3, 1, 4/3, 4/3, 1 y 2/3;
% arranque lento de 0.5 s). Lecturas obligadas: en el corredor con la mezcla el modelo no le gana a la regla reactiva;
% en la red reproduce el ciclo minimo (mix3 y mixa3) y no le gana a la regla simple. veredicto3_comparaciones.csv:
% mix1 frente al actuado ajustado +7.07 y +6.73 (FALLA) y mixa1 +8.33 y +9.04; frente a la regla reactiva +0.31 y
% +0.76; per1 frente al actuado ajustado +1.95 y +1.39 (EMPATA) y frente al fijo por bloques -12.36 y -11.73

\subsection*{Por qué el modelo no supera al fijo en la red con rotonda ni en la malla con demanda nominal}

Esta lectura es posterior al cierre, con los datos del estudio, y no
cambia ningún veredicto; la Sección~\ref{sec:donde-gana} la resume.

\textbf{En esos tres casos ningún control supera al fijo ajustado al
segundo.} La Tabla~\ref{tab:plan3-margen} compara, en cada caso, la
mejor de las otras nueve referencias con el modelo citado. En P1, P2 y
P6 hay margen: el actuado ajustado ya mejora al fijo, y el modelo lo
mejora más. En P3, P4 y P5 la mejor referencia solo iguala al fijo.

\begin{table}[H]
\setstretch{1}
\centering
\footnotesize
\setlength{\tabcolsep}{3pt}
% Generado por tablas_plan3.py a partir de plan3/salidas/. No editar a mano.
% Segundo estudio, llegadas aleatorias, tanda 1: retraso total con el fijo ajustado al segundo, la mejor
% de las otras referencias frente a el y el modelo citado. Lectura posterior al cierre (descriptiva).
% fuente: plan3/salidas/veredicto3_comparaciones.csv (tanda T1, referencia 'fijo fino', claves ref_*:
% media_ref y el menor cambio) y veredicto3_celdas.csv (cambio_T1 del modelo citado).
% Entorno: \setstretch{1}, \footnotesize y \setlength{\tabcolsep}{3pt}.
\begin{tabular}{>{\raggedright\arraybackslash}p{4.6cm}>{\raggedright\arraybackslash}p{2.2cm}>{\raggedright\arraybackslash}p{4.2cm}>{\raggedright\arraybackslash}p{2.0cm}}
\toprule
Caso & Retraso con el fijo ajustado al segundo (s/veh) & La mejor de las otras referencias (\%) & Modelo citado (\%) \\
\midrule
P1, corredor, demanda nominal (1810~veh/h) & 29.2 & $-13.5$ (actuado ajustado) & $-16.3$ \\
P2, corredor, demanda alta (3620~veh/h) & 458.4 & $-1.0$ (actuado ajustado) & $-19.8$ \\
P3, red con rotonda, demanda nominal (2150~veh/h) & 22.7 & $0.0$ (actuado ajustado y 2 más) & $+1.6$ \\
P4, red con rotonda, demanda alta (2580~veh/h) & 28.7 & $0.0$ (fijo por franja) & $-0.5$ \\
P5, malla, demanda nominal (4350~veh/h) & 22.5 & $0.0$ (fijo por franja) & $+0.2$ \\
P6, malla, demanda alta (8700~veh/h) & 51.3 & $-16.0$ (actuado ajustado) & $-25.7$ \\
\bottomrule
\end{tabular}

\caption[Segundo estudio: margen de mejora en cada caso]{Margen de mejora con llegadas aleatorias, en la tanda 1: retraso total con el fijo ajustado al segundo, cambio de la mejor de las otras nueve referencias frente a él (entre paréntesis, cuál es) y cambio del modelo citado.}
\label{tab:plan3-margen}
\end{table}
% fuente: plan3/salidas/veredicto3_comparaciones.csv (tanda T1, referencia "fijo fino", claves ref_*) y
% veredicto3_celdas.csv (cambio_T1)

\textbf{En la red, la diferencia es un segundo de desfase.} El fijo
ajustado al segundo de P3 es el ciclo mínimo con el semáforo del medio
1~s después que los otros dos, y ese segundo explica el 1.6\,\% que
separa al modelo del fijo en la evaluación. En el ajuste, el retraso
baja de 22.31 a 22.01~s/veh al correr ese semáforo 1~s; con 2~s vuelve
a 22.31, y con 1~s antes sube a 22.47. El modelo cambia de fase apenas
se cumple el verde mínimo, en los tres semáforos a la vez, y decide
cada 5~s (cada 2~s el modelo completo): no puede producir un desfase de
1~s. En P4, el fijo ajustado al segundo es el ciclo mínimo con dos
semáforos corridos 1~s.

\textbf{En la malla con demanda nominal, el retraso casi no depende del
control.} El ciclo mínimo, el fijo ajustado desde 10~s y la regla
simple quedan a un 0.2\,\% del fijo ajustado al segundo; el fijo por
franja lo iguala, y el modelo queda a un 0.2 y un 0.3\,\%.
% fuente: plan3/poi_red_r10f.add.xml y plan3/poi_red_trivial.add.xml (mismas fases, 10/3/10/3 s; solo cambia el
% offset de semaforo_C2, 1 frente a 0) y plan3/poi_red_alta_r10f.add.xml (offsets 0, 1 y 25 con ciclo de 26 s);
% plan3/barridos/poi_red/desfases.csv, semillas 2001-2010 (retraso_total_prom de 0/0/0, 0/1/0, 0/2/0 y 0/25/0: 22.31,
% 22.01, 22.31 y 22.47); veredicto3_comparaciones.csv (P3, T1: ciclo minimo y modelo, 23.085 s/veh; fijo ajustado al
% segundo, 22.730; P5, T1: ciclo minimo, fijo entregado y regla simple +0.16 %, fijo por franja 0.00 %)

\subsection*{Lo esperado y lo obtenido}

El registro incluye lo que se esperaba, para contrastarlo con el
resultado (Tabla~\ref{tab:plan3-esperado}). Dos resultados no estaban previstos:
decidir cada 2~s no ayuda por sí solo, y el modelo completo, que no
tenía predicción, es el que más mejora.

\begin{table}[H]
\setstretch{1}
\centering
\footnotesize
\setlength{\tabcolsep}{3pt}
\begin{tabular}{>{\raggedright\arraybackslash}p{6.6cm}>{\raggedright\arraybackslash}p{7.0cm}}
\toprule
Qué se esperaba & Qué se obtuvo \\
\midrule
Llegadas regulares: empate en P6 frente al fijo ajustado al segundo & EMPATA \\
\specialrule{0.25pt}{1.5pt}{1.5pt}
Llegadas regulares: un resumen en un número más cerca de cero que frente al fijo ajustado desde 10~s & $-0.4$\,\% en las dos tandas \\
\specialrule{0.25pt}{1.5pt}{1.5pt}
Llegadas aleatorias: P2 y P6, candidatos a FUNCIONA & FUNCIONA en P2 y P6 (en P2, sin estabilizar) \\
\specialrule{0.25pt}{1.5pt}{1.5pt}
Llegadas aleatorias: en P1, MEJORA o menos, porque el actuado también mejora al fijo & IPPO con mensajes corregido, MEJORA; el modelo completo, FUNCIONA \\
\specialrule{0.25pt}{1.5pt}{1.5pt}
Decidir cada 2~s ayudaría con llegadas aleatorias & No pasa en ninguno de los 3 casos \\
\specialrule{0.25pt}{1.5pt}{1.5pt}
Entrenar 400 episodios mejoraría entre un 1 y un 3\,\% en el corredor & En el corredor no pasa; en P6 mejora un 5\,\% \\
\specialrule{0.25pt}{1.5pt}{1.5pt}
Modelo completo: sin predicción, con riesgo de romper la coordinación & Mejora a IPPO con mensajes corregido en 4 de 6 casos \\
\bottomrule
\end{tabular}
\caption[Segundo estudio: lo esperado y lo obtenido]{Lo que el registro previo esperaba y lo que se obtuvo.}
\label{tab:plan3-esperado}
\end{table}
% fuente: plan3/PRERREGISTRO.md §15 (prediccion) y plan3/salidas/veredicto3.md (veredictos, hipotesis H9, H14 y H15)

\subsection*{Desviaciones del registro}

\begin{itemize}
  \item \textbf{Semillas de evaluación fuera de la evaluación a
    ciegas.} Una revisión del código simuló por error tres semillas de
    evaluación con llegadas regulares, con controles del primer estudio
    y antes de que existiera ningún modelo del estudio. Esos archivos se
    borraron sin abrir, y el código rechaza esas semillas fuera de la
    evaluación a ciegas.
  \item \textbf{Código de las referencias.} Los programas que ajustan
    las referencias tienen correcciones posteriores al registro y
    anteriores a todo ajuste; la versión usada quedó anotada al
    congelar las referencias.
  \item \textbf{Entrenamientos sin estabilizar.} En P2 con llegadas
    aleatorias, en P5 con regulares y en el corredor con perfil de
    demanda, el entrenamiento de la versión evaluada no se estabilizó;
    su veredicto se calcula igual y lleva esa marca. Como en el primer
    estudio, no se agregaron episodios. Con llegadas aleatorias se
    estabilizan 41 de los 48 entrenamientos de IPPO con mensajes
    corregido y del modelo completo (cinco y tres por caso): 25 de 30 y
    16 de 18.
    % fuente: plan3/salidas/veredicto3_bases.csv, filas con llegadas aleatorias de las familias v2poi (IPPO con
    % mensajes corregido, 5 bases x 6 casos) y p3completopoi (modelo completo, 3 x 6), columna convergio: 41 "si" de 48
    % (25 de 30 y 16 de 18); sin estabilizar: P2 v2poi 2042, P2 completo 42 y 4042 (la version evaluada), P5 v2poi 42 y
    % 14042, P6 v2poi 2042 y 16042
  \item \textbf{Vehículos sin entrar.} En P2 y P6, con el ciclo mínimo,
    quedan vehículos sin entrar a la red al cerrar la simulación: su
    retraso está subestimado y la ventaja del modelo frente a él es
    conservadora.
  \item \textbf{Entrenamientos interrumpidos.} Ninguno.
\end{itemize}
% fuente: plan3/doc/desviaciones.md (§1 semillas 8001-8003 en una revision; §2 conteo de cambios de fase; §3 codigo de
% las referencias; §4 y §5 enmiendas); plan3/doc/enmienda_1.md y enmienda_2.md; plan3/salidas/veredicto3.md
% ("PENDIENTE", "Vehiculos sin entrar al cierre", "Entrenamientos caidos o sin publicar")
% El apartado "Como se reproduce" salio el 15.11.2026: el registro previo, las enmiendas, los programas y las salidas
% estan en plan3/ y sus programas en el Anexo B (tab:programas-segundo y tab:programas-lima).
